\chapter{2023年新高考I卷}

\section{单选题}

\begin{problem}
已知集合 \(M=\{-2,-1,0,1,2\}\)，\(N=\{x\mid x^2-x-6\ge 0\}\)，则 \(M\cap N=\)（\quad）

\choices
    {\(\{-2,-1,0,1\}\)}
    {\(\{0,1,2\}\)}
    {\(\{-2\}\)}
    {\(\{2\}\)}
\end{problem}

\begin{answer}
C
\end{answer}

\begin{solution}
由
\(x^2-x-6\ge 0 \iff (x+2)(x-3)\ge 0 \iff x\le -2 \text{\ 或\ } x\ge 3\)，
得 \(N=(-\infty,-2]\cup[3,+\infty)\). 因此
\(M\cap N=\{-2\}\).
故选 C.
\end{solution}

\begin{problem}
已知 \(z=\frac{1-\i}{2+2\i}\)，则 \(z-\overline z=\)（\quad）

\choices
    {\(-\i\)}
    {\(\i\)}
    {\(0\)}
    {\(1\)}
\end{problem}

\begin{answer}
A
\end{answer}

\begin{solution}
化简得
\(z=\frac{(1-\i)(2-2\i)}{(2+2\i)(2-2\i)} =\frac{-4\i}{8} =-\frac{\i}{2}\)，
所以 \(\overline z=\frac{\i}{2}\). 故
\(z-\overline z=-\frac{\i}{2}-\frac{\i}{2}=-\i\).
故选 A.
\end{solution}

\begin{problem}
已知向量 \(\bs{a}=(1,1)\)，\(\bs{b}=(1,-1)\)，若 \((\bs{a}+\lambda\bs{b})\perp(\bs{a}+\mu\bs{b})\)，则（\quad）

\choices
    {\(\lambda+\mu=1\)}
    {\(\lambda+\mu=-1\)}
    {\(\lambda\mu=1\)}
    {\(\lambda\mu=-1\)}
\end{problem}

\begin{answer}
D
\end{answer}

\begin{solution}
由垂直得
\((\bs{a}+\lambda\bs{b})\cdot(\bs{a}+\mu\bs{b})=0\).
又
\(\bs{a}\cdot\bs{a}=2\)，\(\bs{b}\cdot\bs{b}=2\)，\(\bs{a}\cdot\bs{b}=0\)，
所以
\(2+2\lambda\mu=0\)，
故 \(\lambda\mu=-1\). 故选 D.
\end{solution}

\begin{problem}
设函数 \(f(x)=2^{x(x-a)}\) 在区间 \((0,1)\) 单调递减，则 \(a\) 的取值范围是（\quad）

\choices
    {\((-\infty,-2]\)}
    {\([-2,0)\)}
    {\((0,2]\)}
    {\([2,+\infty)\)}
\end{problem}

\begin{answer}
D
\end{answer}

\begin{solution}
因为 \(y=2^u\) 在 \(\R\) 上单调递增，所以只需
\(u=x(x-a)=x^2-ax\)
在 \((0,1)\) 上单调递减. 其对称轴为
\(x=\frac{a}{2}\).
要使抛物线在 \((0,1)\) 上整体递减，应有 \(\frac{a}{2}\ge 1\). 故
\(a\ge 2\).

\solutionfigure{\bitmapfigure[max width=0.42\linewidth]{img/2023/new_gaokao_paper_1/q4_solution_fig1.png}}

故选 D.
\end{solution}

\begin{problem}
设椭圆 \(C_1:\frac{x^2}{a^2}+y^2=1\ (a>1)\)，\(C_2:\frac{x^2}{4}+y^2=1\) 的离心率分别为 \(e_1\)，\(e_2\)，若 \(e_2=\sqrt3\,e_1\)，则 \(a=\)（\quad）

\choices
    {\(\frac{2\sqrt3}{3}\)}
    {\(\sqrt2\)}
    {\(\sqrt3\)}
    {\(\sqrt6\)}
\end{problem}

\begin{answer}
A
\end{answer}

\begin{solution}
由椭圆离心率公式，
\(e_1=\frac{\sqrt{a^2-1}}{a}\)，\(e_2=\frac{\sqrt{4-1}}{2}=\frac{\sqrt3}{2}\).
由 \(e_2=\sqrt3\,e_1\) 得
\(\frac{\sqrt3}{2}=\sqrt3\cdot\frac{\sqrt{a^2-1}}{a}\)，
即
\(\frac12=\frac{\sqrt{a^2-1}}{a}\).
解得
\(a=\frac{2\sqrt3}{3}\).
故选 A.
\end{solution}

\begin{problem}
过点 \((0,-2)\) 与圆 \(x^2+y^2-4x-1=0\) 相切的两直线的夹角为 \(\alpha\)，则 \(\sin\alpha=\)（\quad）

\choices
    {\(1\)}
    {\(\frac{\sqrt{15}}{4}\)}
    {\(\frac{\sqrt{10}}{4}\)}
    {\(\frac{\sqrt6}{4}\)}
\end{problem}

\begin{answer}
B
\end{answer}

\begin{solution}
圆可化为
\((x-2)^2+y^2=5\)，
故圆心 \(C(2,0)\)，半径 \(r=\sqrt5\). 记 \(P(0,-2)\)，切点为 \(A\)，\(B\). 由图可知
\(\sin\alpha=\sin\angle APB\).
在直角三角形 \(PAC\) 中，
\(PC=2\sqrt2\)，\(AC=\sqrt5\)，\(PA=\sqrt{PC^2-AC^2}=\sqrt3\).
所以
\(\cos\angle APC=\frac{\sqrt3}{2\sqrt2}\)，\(\sin\angle APC=\frac{\sqrt5}{2\sqrt2}\).
又 \(\angle APB=2\angle APC\)，故
\[
\sin\alpha=\sin 2\angle APC=2\sin\angle APC\cos\angle APC=\frac{\sqrt{15}}{4}
\]

\solutionfigure{\bitmapfigure[max width=0.42\linewidth]{img/2023/new_gaokao_paper_1/q6_solution_fig1.png}}

故选 B.
\end{solution}

\begin{problem}
记 \(S_n\) 为数列 \(\{a_n\}\) 的前 \(n\) 项和. 设甲：\(\{a_n\}\) 为等差数列，乙：\(\left\{\frac{S_n}{n}\right\}\) 为等差数列，则（\quad）

\choices
    {甲是乙的充分条件但不是必要条件}
    {甲是乙的必要条件但不是充分条件}
    {甲是乙的充要条件}
    {甲既不是乙的充分条件也不是乙的必要条件}
\end{problem}

\begin{answer}
C
\end{answer}

\begin{solution}
若 \(\{a_n\}\) 为等差数列，则
\(S_n=An^2+Bn\)，
从而
\(\frac{S_n}{n}=An+B\)，
故乙成立.

反过来，若 \(\left\{\frac{S_n}{n}\right\}\) 为等差数列，可设
\(\frac{S_n}{n}=pn+q\)，
则
\(S_n=pn^2+qn\)，
满足等差数列前 \(n\) 项和的形式特征，因此 \(\{a_n\}\) 为等差数列. 故甲，乙互为充要条件. 故选 C.
\end{solution}

\begin{problem}
已知 \(\sin(\alpha-\beta)=\frac13\)，\(\cos\alpha\sin\beta=\frac16\)，则 \(\cos(2\alpha+2\beta)=\)（\quad）

\choices
    {\(\frac79\)}
    {\(\frac19\)}
    {\(-\frac19\)}
    {\(-\frac79\)}
\end{problem}

\begin{answer}
B
\end{answer}

\begin{solution}
由
\(\sin(\alpha-\beta)=\sin\alpha\cos\beta-\cos\alpha\sin\beta=\frac13\)
及 \(\cos\alpha\sin\beta=\frac16\)，得
\(\sin\alpha\cos\beta=\frac12\).
所以
\(\sin(\alpha+\beta)=\sin\alpha\cos\beta+\cos\alpha\sin\beta=\frac12+\frac16=\frac23\).
故
\(\cos(2\alpha+2\beta)=1-2\sin^2(\alpha+\beta)=1-2\cdot\left(\frac23\right)^2=\frac19\).
故选 B.
\end{solution}
\section{多选题}

\begin{problem}
有一组样本数据 \(x_1\)，\(x_2\)，\(\dots\)，\(x_6\)，其中 \(x_1\) 是最小值，\(x_6\) 是最大值，则（\quad）

\choices
    {\(x_2,x_3,x_4,x_5\) 的平均数等于 \(x_1,x_2,\dots,x_6\) 的平均数}
    {\(x_2\)，\(x_3\)，\(x_4\)，\(x_5\) 的中位数等于 \(x_1\)，\(x_2\)，\(\dots\)，\(x_6\) 的中位数}
    {\(x_2\)，\(x_3\)，\(x_4\)，\(x_5\) 的标准差不小于 \(x_1\)，\(x_2\)，\(\dots\)，\(x_6\) 的标准差}
    {\(x_2\)，\(x_3\)，\(x_4\)，\(x_5\) 的极差不大于 \(x_1\)，\(x_2\)，\(\dots\)，\(x_6\) 的极差}
\end{problem}

\begin{answer}
BD
\end{answer}

\begin{solution}
A 不一定成立，去掉最小值和最大值后平均数可能改变.

设数据按从小到大排列，则 \(6\) 个数据的中位数与中间 \(4\) 个数据的中位数都为
\(\frac{x_3+x_4}{2}\)，
故 B 正确.

C 不一定成立，去掉离平均数较远的两端点后，标准差可能减小.

又
\((x_6-x_1)-(x_5-x_2)=(x_6-x_5)+(x_2-x_1)\ge 0\)，
故
\(x_5-x_2\le x_6-x_1\).
所以 D 正确. 故选 BD.
\end{solution}

\begin{problem}
噪声污染问题越来越受到重视. 用声压级来度量噪声的强度，定义声压级 \(L_P=20\lg\frac{p}{p_0}\)，其中常数 \(p_0(p_0>0)\) 是听觉下限阈值，\(p\) 是实际声压. 下表为不同声源的声压级：

\begin{center}
\begin{tabular}{|c|c|c|}
\hline
声源 & 与声源的距离/m & 声压级/dB \\
\hline
燃油汽车 & 10 & 60\textasciitilde90 \\
\hline
混合动力汽车 & 10 & 50\textasciitilde60 \\
\hline
电动汽车 & 10 & 40 \\
\hline
\end{tabular}
\end{center}

已知在距离燃油汽车，混合动力汽车，电动汽车 \(10\text{m}\) 处测得实际声压分别为 \(p_1\)，\(p_2\)，\(p_3\)，则（\quad）

\choices
    {\(p_1\ge p_2\)}
    {\(p_2>10p_3\)}
    {\(p_3=100p_0\)}
    {\(p_1\le 100p_2\)}
\end{problem}

\begin{answer}
ACD
\end{answer}

\begin{solution}
由
\(L_P=20\lg\frac{p}{p_0}\)
得
\(p=p_0\cdot 10^{L_P/20}\).
因此声压级越大，实际声压越大，A 正确.

对电动汽车，有
\(p_3=p_0\cdot 10^{40/20}=100p_0\)，
故 C 正确.

对混合动力汽车，
\(p_0\cdot 10^{50/20}\le p_2\le p_0\cdot 10^{60/20}\)，
即
\(100\sqrt{10}\,p_0\le p_2\le 1000p_0\)，
故 \(p_2\le 10p_3\)，所以 B 错误.

对燃油汽车，
\(p_0\cdot 10^{60/20}\le p_1\le p_0\cdot 10^{90/20}\)，
即
\(1000p_0\le p_1\le 10000\sqrt{10}\,p_0\).
而 \(100p_2\) 的范围覆盖该区间，故 \(p_1\le 100p_2\)，D 正确. 故选 ACD.
\end{solution}

\begin{problem}
已知函数 \(f(x)\) 的定义域为 \(\R\)，且 \(f(xy)=y^2f(x)+x^2f(y)\)，则（\quad）

\choices
    {\(f(0)=0\)}
    {\(f(1)=0\)}
    {\(f(x)\) 是偶函数}
    {\(x=0\) 为 \(f(x)\) 的极小值点}
\end{problem}

\begin{answer}
ABC
\end{answer}

\begin{solution}
令 \(x=y=0\)，得
\(f(0)=0\).
令 \(x=y=1\)，得
\(f(1)=2f(1)\)，
故 \(f(1)=0\).

令 \(y=-1\)，得
\(f(-x)=f(x)+x^2f(-1)\).
再令 \(x=y=-1\)，得
\(f(1)=2f(-1)\)，
由 \(f(1)=0\) 得 \(f(-1)=0\)，从而
\(f(-x)=f(x)\).
故 \(f\) 为偶函数.

D 不一定成立，例如零函数满足题意，但 \(x=0\) 不是它的极小值点. 故选 ABC.
\end{solution}

\begin{problem}
下列物体中，能够被整体放入棱长为 \(1\)（单位：m）的正方体容器（容器壁厚度忽略不计）内的有（\quad）

\choices
    {直径为 \(0.99\text{m}\) 的球体}
    {所有棱长均为 \(1.4\text{m}\) 的正四面体}
    {底面直径为 \(0.01\text{m}\)，高为 \(1.8\text{m}\) 的圆柱体}
    {底面直径为 \(1.2\text{m}\)，高为 \(0.01\text{m}\) 的圆柱体}
\end{problem}

\begin{answer}
ABD
\end{answer}

\begin{solution}
A 项，因为正方体的内切球直径为 \(1\text{m}\)，所以直径为 \(0.99\text{m}\) 的球体可以放入正方体容器，故 A 项正确；

B 项，看到正方体和正四面体，要想到由正方体的面对角线可以构成正四面体，如图，蓝色正四面体的棱长为 \(\sqrt2\)，比 \(1.4\) 大，从而所有棱长均为 \(1.4\text{m}\) 的正四面体可以放入正方体容器，故 B 项正确；

C 项，注意到底面直径很小，圆柱很细长，不妨将其近似成线段，故先看 \(1.8\text{m}\) 的线段能否放入正方体. 如图，正方体的棱长为 \(1\)，则正方体表面上任意两点之间距离的最大值为 \(\sqrt3<1.8\)，所以高为 \(1.8\text{m}\) 的圆柱体不可能放入该正方体，故 C 项错误；

D 项，注意到圆柱的高很小，不妨将圆柱近似看成圆，故先分析直径为 \(1.2\text{m}\) 的圆能否放入正方体. 为了研究这一问题，我们先找正方体尽可能大的截面. 如图，\(E\)，\(F\)，\(G\)，\(H\)，\(I\)，\(J\) 分别为所在棱的中点，则 \(EFGHIJ\) 是边长为 \(\frac{\sqrt2}{2}\) 的正六边形，其内切圆如图，其中 \(K\) 为 \(HI\) 中点，则内切圆半径
\(r=OK=\frac{\sqrt2}{2}\times\frac{\sqrt3}{2}=\frac{\sqrt6}{4}\)，
直径
\(2r=\frac{\sqrt6}{2}>1.2\)，
所以可以想象，底面直径为 \(1.2\text{m}\)，高为 \(0.01\text{m}\) 的圆柱体能够放进正方体容器，故 D 项正确.

\solutionfigure{\bitmapfigure[max width=0.62\linewidth]{img/2023/new_gaokao_paper_1/q12_solution_fig1.png}}

故选 ABD.
\end{solution}
\section{填空题}

\begin{problem}
某学校开设了 \(4\) 门体育类选修课和 \(4\) 门艺术类选修课. 学生需从这 \(8\) 门课中选 \(2\) 门或 \(3\) 门课，并且每类选修课至少选 \(1\) 门，则不同的选课方案共有\fillinblank{}种. （用数字作答）
\end{problem}

\begin{answer}
64
\end{answer}

\begin{solution}
若选 \(2\) 门，只能体育类，艺术类各选 \(1\) 门，共
\(\binom41\binom41=16\)
种.

若选 \(3\) 门，可分为“体 \(1\) 艺 \(2\)”或“体 \(2\) 艺 \(1\)”两类，共
\(\binom41\binom42+\binom42\binom41=48\)
种.

故总数为
\(16+48=64\).
\end{solution}

\begin{problem}
在正四棱台 \(ABCD-A_1B_1C_1D_1\) 中，\(AB=2\)，\(A_1B_1=1\)，\(AA_1=\sqrt2\)，则该棱台的体积为\fillinblank{}.
\end{problem}

\begin{answer}
\(\dfrac{7\sqrt6}{6}\)
\end{answer}

\begin{solution}
体积公式为
\(V=\frac13\bigl(S+S'+\sqrt{SS'}\bigr)h\)，
故只需求高 \(h\).
在截面 \(AA_1C_1C\) 中，作 \(A_1E\perp AC\)，\(C_1F\perp AC\). 则
\(EF=A_1C_1=\sqrt2\)，\(AC=2\sqrt2\).
所以
\(AE=\frac{AC-EF}{2}=\frac{\sqrt2}{2}\).
由 \(AA_1=\sqrt2\) 得
\[
h=A_1E=\sqrt{AA_1^2-AE^2}=\sqrt{2-\frac12}=\sqrt{\frac32}=\frac{\sqrt6}{2}
\]

\solutionfigure{\bitmapfigure[max width=0.5\linewidth]{img/2023/new_gaokao_paper_1/q14_solution_fig1.png}}

又上，下底面积分别为
\(S'=1\)，\(S=4\)，
故
\(V=\frac13(4+1+2)\cdot\frac{\sqrt6}{2} =\frac{7\sqrt6}{6}\).
\end{solution}

\begin{problem}
已知函数 \(f(x)=\cos\omega x-1\)（\(\omega>0\)） 在区间 \([0,2\pi]\) 上有且仅有 \(3\) 个零点，则 \(\omega\) 的取值范围是\fillinblank{}.
\end{problem}

\begin{answer}
\([2,3)\)
\end{answer}

\begin{solution}
由 \(f(x)=0\) 得
\(\cos\omega x=1\).
所以问题等价于 \(y=\cos\omega x\) 在 \([0,2\pi]\) 上恰有 \(3\) 个最大值点. 由图可知应有
\(\frac{4\pi}{\omega}\le 2\pi<\frac{6\pi}{\omega}\).
解得
\(2\le \omega<3\).

\solutionfigure{\bitmapfigure[max width=0.6\linewidth]{img/2023/new_gaokao_paper_1/q15_solution_fig1.png}}
\end{solution}

\begin{problem}
已知双曲线 \(C:\frac{x^2}{a^2}-\frac{y^2}{b^2}=1\)（\(a>0\)，\(b>0\)） 的左，右焦点分别为 \(F_1\)，\(F_2\). 点 \(A\) 在 \(C\) 上，点 \(B\) 在 \(y\) 轴上，\(F_1A\perp F_1B\)，\(\overrightarrow{F_2A}=-\frac23\overrightarrow{F_2B}\)，则 \(C\) 的离心率为\fillinblank{}.
\end{problem}

\begin{answer}
\(\dfrac{3\sqrt5}{5}\)
\end{answer}

\begin{solution}
设 \(AF_2=2m\). 由
\(\overrightarrow{F_2A}=-\frac23\overrightarrow{F_2B}\)
可得 \(BF_2=3m\)，于是
\(AB=AF_2+BF_2=5m\).
又由对称性 \(BF_1=BF_2=3m\). 因为 \(F_1A\perp F_1B\)，故
\(AF_1=\sqrt{AB^2-BF_1^2}=4m\).
点 \(A\) 在双曲线右支上，所以
\(AF_1-AF_2=2a\)，
即
\(4m-2m=2a\)，
故 \(m=a\).
在 \(\triangle BF_1F_2\) 中，
\(\cos\angle F_1BF_2=\frac{BF_1^2+BF_2^2-F_1F_2^2}{2BF_1\cdot BF_2} =\frac{9a^2-2c^2}{9a^2}\).
而
\(\cos\angle ABF_1=\frac{BF_1}{AB}=\frac{3}{5}\).
由图知 \(\angle ABF_1=\angle F_1BF_2\)，故
\(\frac{9a^2-2c^2}{9a^2}=\frac35\).
解得
\(\frac{c}{a}=\frac{3\sqrt5}{5}\).

\solutionfigure{\bitmapfigure[max width=0.5\linewidth]{img/2023/new_gaokao_paper_1/q16_solution_fig1.png}}
\end{solution}
\section{解答题}

\begin{problem}
已知在 \(\triangle ABC\) 中，\(A+B=3C\)，\(2\sin(A-C)=\sin B\).
\begin{enumerate}
\item 求 \(\sin A\)；
\item 设 \(AB=5\)，求 \(AB\) 边上的高.
\end{enumerate}
\end{problem}

\begin{solution}
\begin{enumerate}
\item 因为
\(A+B=\pi-C=3C\)，
故
\(C=\frac{\pi}{4}\)，\(B=\frac{3\pi}{4}-A\).
要求的是 \(\sin A\)，故用 \(C=\frac{\pi}{4}\) 和 \(A+B=\frac{3\pi}{4}\) 将 \(2\sin(A-C)=\sin B\) 消元，把变量统一成 \(A\).
代入 \(2\sin(A-C)=\sin B\)，得
\(2\sin\left(A-\frac{\pi}{4}\right)=\sin\left(\frac{3\pi}{4}-A\right)\).
所以
\(2\left(\sin A\cos\frac{\pi}{4}-\cos A\sin\frac{\pi}{4}\right)=\sin\frac{3\pi}{4}\cos A-\cos\frac{3\pi}{4}\sin A\)，
整理得
\(\cos A=\frac13\sin A\).
由
\(\sin^2A+\cos^2A=1\)
得
\(\sin^2A+\frac19\sin^2A=1\)，
故
\(\sin A=\pm\frac{3\sqrt{10}}{10}\).
结合 \(0<A<\pi\) 可得
\(\sin A=\frac{3\sqrt{10}}{10}\).

\item 设内角 \(A\)，\(B\)，\(C\) 的对边分别为 \(a\)，\(b\)，\(c\)，则 \(c=AB=5\). 如图，\(AB\) 边上的高 \(CD=a\sin B\).

已知 \(A\)，\(C\)，故 \(\sin B\) 可用内角和为 \(\pi\) 来求：
\(\sin B=\sin\left(\frac{3\pi}{4}-A\right)=\frac{\sqrt2}{2}\cos A+\frac{\sqrt2}{2}\sin A=\frac{2}{\sqrt5}\).
再求 \(a\)，已知条件有 \(C\)，\(c\)，\(\sin A\)，故用正弦定理求 \(a\). 由正弦定理，
\(\frac{a}{\sin A}=\frac{c}{\sin C}\)，
得
\(a=\frac{c\sin A}{\sin C}=\frac{5\cdot \frac{3\sqrt{10}}{10}}{\frac{\sqrt2}{2}}=3\sqrt5\).
因此
\(CD=a\sin B=3\sqrt5\cdot \frac{2}{\sqrt5}=6\).

\solutionfigure{\bitmapfigure[max width=0.42\linewidth]{img/2023/new_gaokao_paper_1/q17_solution_fig1.png}}
\end{enumerate}
\end{solution}

\begin{problem}
如图，在正四棱柱 \(ABCD-A_1B_1C_1D_1\) 中，\(AB=2\)，\(AA_1=4\). 点 \(A_2\)，\(B_2\)，\(C_2\)，\(D_2\) 分别在棱 \(AA_1\)，\(BB_1\)，\(CC_1\)，\(DD_1\) 上，且 \(AA_2=1\)，\(BB_2=DD_2=2\)，\(CC_2=3\).

\begin{enumerate}
\item 证明：\(B_2C_2\parallel A_2D_2\)；
\item 点 \(P\) 在棱 \(BB_1\) 上，当二面角 \(P-A_2C_2-D_2\) 为 \(150^\circ\) 时，求 \(B_2P\).
\end{enumerate}

\bitmapfigure[max width=0.42\linewidth]{img/2023/new_gaokao_paper_1/q18_fig1.png}
\end{problem}

\begin{solution}
\begin{enumerate}
\item 以 C 为原点建立空间直角坐标系，如图可取
\(B_2(0,2,2)\)，\(C_2(0,0,3)\)，\(A_2(2,2,1)\)，\(D_2(2,0,2)\).
于是
\(\overrightarrow{B_2C_2}=(0,-2,1)\)，\(\overrightarrow{A_2D_2}=(0,-2,1)\).
故 \(B_2C_2\parallel A_2D_2\).

\item 设
\(P(0,2,a)\)（\(0\le a\le 4\)）.
则平面 \(PA_2C_2\) 的一个法向量可取
\(\mathbf{m}=(1-a,a-3,-2)\)，
平面 \(A_2C_2D_2\) 的一个法向量可取
\(\mathbf{n}=(1,1,2)\).
因为二面角 \(P-A_2C_2-D_2\) 为 \(150^\circ\)，故法向量夹角的锐角为 \(30^\circ\)，从而
\(\frac{|\mathbf{m}\cdot \mathbf{n}|}{|\mathbf{m}||\mathbf{n}|}=\frac{\sqrt3}{2}\).
而
\(\mathbf{m}\cdot \mathbf{n}=(1-a)+(a-3)-4=-6\)，
且
\(|\mathbf{n}|=\sqrt6\).
于是
\(\frac{6}{\sqrt{(1-a)^2+(a-3)^2+4}\cdot \sqrt6}=\frac{\sqrt3}{2}\).
解得
\(a=1\text{\ 或\ }3\).
所以
\(B_2P=|a-2|=1\).

\solutionfigure{\bitmapfigure[max width=0.42\linewidth]{img/2023/new_gaokao_paper_1/q18_solution_fig1.png}}
\end{enumerate}
\end{solution}

\begin{problem}
已知函数 \(f(x)=a(\e^x+a)-x\).
\begin{enumerate}
\item 讨论 \(f(x)\) 的单调性；
\item 证明：当 \(a>0\) 时，\(f(x)>2\ln a+\frac32\).
\end{enumerate}
\end{problem}

\begin{solution}
\begin{enumerate}
\item
\(f'(x)=a\e^x-1\).
当 \(a\le 0\) 时，\(f'(x)<0\)，故 \(f\) 在 \(\R\) 上单调递减.

当 \(a>0\) 时，
\(f'(x)=0 \iff x=\ln\frac1a\).
因此 \(f\) 在 \((-\infty,\ln\frac1a)\) 上单调递减，在 \((\ln\frac1a,+\infty)\) 上单调递增.

\item 由上一项知，当 \(a>0\) 时，\(f\) 的最小值为
\(f\left(\ln\frac1a\right) =a\left(\frac1a+a\right)-\ln\frac1a =1+a^2+\ln a\).
故只需证明
\(1+a^2+\ln a>2\ln a+\frac32\)，
即
\(a^2-\ln a-\frac12>0\).
令
\(g(a)=a^2-\ln a-\frac12\)（\(a>0\)），
则
\(g'(a)=2a-\frac1a=\frac{2a^2-1}{a}\).
所以 \(g(a)\) 在 \((0,\frac{\sqrt2}{2})\) 上单调递减，在 \((\frac{\sqrt2}{2},+\infty)\) 上单调递增.
故
\(g(a)\ge g\left(\frac{\sqrt2}{2}\right) =\frac12-\ln\frac{\sqrt2}{2}-\frac12 =-\ln\frac{\sqrt2}{2}>0\).
所以
\(f(x)\ge f\left(\ln\frac1a\right)>2\ln a+\frac32\).
\end{enumerate}
\end{solution}

\begin{problem}
设等差数列 \(\{a_n\}\) 的公差为 \(d\)，且 \(d>1\). 令 \(b_n=\frac{n^2+n}{a_n}\)，记 \(S_n\)，\(T_n\) 分别为数列 \(\{a_n\}\)，\(\{b_n\}\) 的前 \(n\) 项和.
\begin{enumerate}
\item 若 \(3a_2=3a_1+a_3\)，\(S_3+T_3=21\)，求 \(\{a_n\}\) 的通项公式；
\item 若 \(\{b_n\}\) 为等差数列，且 \(S_{99}-T_{99}=99\)，求 \(d\).
\end{enumerate}
\end{problem}

\begin{solution}
\begin{enumerate}
\item 由
\(3a_2=3a_1+a_3\)
得
\(3(a_1+d)=3a_1+(a_1+2d)\)，
故
\(a_1=d\).
又
\(S_3=3a_1+3d\)，\(T_3=\frac{2}{a_1}+\frac{6}{a_1+d}+\frac{12}{a_1+2d}\).
代入 \(a_1=d\) 与 \(S_3+T_3=21\)，可得
\(2d+\frac3d=7\).
解得
\(d=3\text{\ 或\ }\frac12\).
由 \(d>1\)，知 \(d=3\)，于是 \(a_1=3\)，故
\(a_n=a_1+(n-1)d=3n\).

\item 若 \(\{b_n\}\) 为等差数列，则
\(2b_2=b_1+b_3\).
代入
\(b_1=\frac{2}{a_1}\)，\(b_2=\frac{6}{a_1+d}\)，\(b_3=\frac{12}{a_1+2d}\)，
化简得
\((a_1-d)(a_1-2d)=0\).
故分两种情况：

若 \(a_1=d\)，则 \(a_n=nd\)，从而 \(S_n=\frac{n(n+1)}{2}d\)，\(b_n=\frac{n+1}{d}\)，\(T_n=\frac{n(n+3)}{2d}\).
代入 \(S_{99}-T_{99}=99\)，得
\(99\cdot 50d-\frac{99\cdot 51}{d}=99\).
解得
\(d=\frac{51}{50}\text{\ 或\ }-1\).
由 \(d>1\)，得
\(d=\frac{51}{50}\).

若 \(a_1=2d\)，则
\(a_n=a_1+(n-1)d=(n+1)d\)，从而 \(S_n=\frac{n(a_1+a_n)}{2}=\frac{n(n+3)}{2}d\)，\(b_n=\frac{n^2+n}{(n+1)d}=\frac{n}{d}\)，且
\[
T_n=\frac{n(b_1+b_n)}{2}=\frac{n\left(\frac{1}{d}+\frac{n}{d}\right)}{2}=\frac{n(n+1)}{2d}
\]
代入 \(S_{99}-T_{99}=99\)，得
\(99\cdot 51d-\frac{99\cdot 50}{d}=99\).
解得
\(d=-\frac{50}{51}\text{\ 或\ }1\).
均不满足 \(d>1\)，舍去.

综上，
\(d=\frac{51}{50}\).
\end{enumerate}
\end{solution}

\begin{problem}
甲乙两人投篮，每次由其中一人投篮，规则如下：若命中则此人继续投篮，若未命中则换为对方投篮. 无论之前投篮情况如何，甲每次投篮的命中率均为 \(0.6\)，乙每次投篮的命中率均为 \(0.8\). 由抽签确定第一次投篮的人选，第一次投篮的人是甲，乙的概率各为 \(0.5\).
\begin{enumerate}
\item 求第二次投篮的人是乙的概率；
\item 求第 \(i\) 次投篮的人是甲的概率；
\item 已知：若随机变量 \(X_i\) 服从两点分布，且 \(P(X_i=1)=1-P(X_i=0)=q_i\)，\(\ i=1\)，\(2\)，\(\dots\)，\(n\)，则 \(E\left(\sum_{i=1}^n X_i\right)=\sum_{i=1}^n q_i\).
记前 \(n\) 次（即从第 \(1\) 次到第 \(n\) 次）投篮中甲投篮的次数为 \(Y\)，求 \(E(Y)\).
\end{enumerate}
\end{problem}

\begin{solution}
\begin{enumerate}
\item 记“第 \(i\) 次投篮的人是甲”为事件 \(A_i\). 则
\(P(A_1)=\frac12\).
第二次投篮的人是乙，当且仅当“第一次是甲且甲未中”或“第一次是乙且乙命中”，故
\(P(\text{第2次是乙}) =\frac12(1-0.6)+\frac12\cdot 0.8 =0.6\).

\item 当 \(i\ge 2\) 时，由全概率公式，
\(P(A_i)=P(A_{i-1})\cdot 0.6+[1-P(A_{i-1})]\cdot 0.2 =\frac25P(A_{i-1})+\frac15\).
令
\(u_i=P(A_i)-\frac13\)，
则
\(u_i=\frac25u_{i-1}\).
又
\(u_1=\frac12-\frac13=\frac16\)，
故
\(u_i=\frac16\left(\frac25\right)^{i-1}\).
因此
\(P(A_i)=\frac13+\frac16\left(\frac25\right)^{i-1}\).

\item 设随机变量
\[
X_i=\begin{cases}
1,&\text{第 }i\text{ 次是甲投篮},\\
0,&\text{第 }i\text{ 次是乙投篮}.
\end{cases}
\]
则
\(Y=X_1+X_2+\cdots+X_n\)，
从而
\(E(Y)=\sum_{i=1}^n E(X_i)=\sum_{i=1}^n P(A_i)\).
代入上一项中的结果，
\[
E(Y)=\sum_{i=1}^n\left[\frac13+\frac16\left(\frac25\right)^{i-1}\right]=\frac{n}{3}+\frac16\cdot \frac{1-(\frac25)^n}{1-\frac25}
\]
故
\(E(Y)=\frac{n}{3}+\frac{5}{18}\left[1-\left(\frac25\right)^n\right]\).
\end{enumerate}
\end{solution}

\begin{problem}
在直角坐标系 \(xOy\) 中，点 \(P\) 到 \(x\) 轴的距离等于点 \(P\) 到点 \((0,\frac12)\) 的距离，记动点 \(P\) 的轨迹为 \(W\).
\begin{enumerate}
\item 求 \(W\) 的方程；
\item 已知矩形 \(ABCD\) 有三个顶点在 \(W\) 上，证明：矩形 \(ABCD\) 的周长大于 \(3\sqrt3\).
\end{enumerate}
\end{problem}

\begin{solution}
\begin{enumerate}
\item 由题意得 \(W\) 为抛物线，且准线为 \(y=0\)，焦点为 \(\left(0,\frac12\right)\). 这是由标准抛物线方程 \(x^2=2py\) 向上平移 \(\frac14\) 个单位得到的，故可设为
\(x^2=2py-\frac14\).
因为焦点到准线的距离为 \(p\)，所以 \(p=\frac12\)，故
\(W:\ x^2=y-\frac14\).

\item 不妨设 \(A\)，\(B\)，\(C\) 在抛物线上，且 \(AB\perp BC\)，所以
\(\frac{y_B-y_C}{x_B-x_C}\cdot \frac{y_B-y_A}{x_B-x_A}=-1\)
化简为
\((x_B+x_A)(x_B+x_C)=-1\).
令
\(x_B+x_A=m\)，\(x_B+x_C=-\frac1m\).
由对称性可设 \(|m|\le 1\). 则
\(\frac12\cdot \text{周长}=AB+BC\).
由
\(y=x^2+\frac14\)
及距离公式，可得
\(AB+BC=\sqrt{(y_A-y_B)^2+(x_A-x_B)^2}+\sqrt{(y_C-y_B)^2+(x_C-x_B)^2}\)
\(=|x_A-x_B|\cdot\sqrt{1+m^2}+\frac{|x_C-x_B|}{|m|}\cdot\sqrt{1+m^2}\)
\(\ge \sqrt{1+m^2}\cdot\left(|x_A-x_B|+|x_C-x_B|\right)\)
\(\ge \sqrt{1+m^2}\cdot|x_A-x_C|\)
\(=\sqrt{1+m^2}\cdot\left|m+\frac1m\right| =\sqrt{\frac{(1+m^2)^3}{m^2}}\).
令
\(\phi(t)=\frac{(1+t)^3}{t}\)（\(0<t\le 1\)），
其中 \(t=m^2\). 则
\(\phi'(t)=\frac{(t+1)^2(2t-1)}{t^2}\).
所以 \(\phi(t)\ge \phi\left(\frac12\right)=\frac{27}{4}\). 故
\(\frac12\cdot \text{周长}> \sqrt{\frac{27}{4}}=\frac{3\sqrt3}{2}\).
于是
\(\text{周长}>3\sqrt3\).
\end{enumerate}
\end{solution}
